\documentclass[]{ConspectBook}
\let\enquote\emph

%  1) Charges
%  2) ChargeSystem
%  3) EBMotion
%  4) ElEnergy
%  5) ElMagInduction
%  6) EMOscillations
%  7) EMWaves
%  8) FileldInDielectrics
%  9) FileldInMetall
% 10) MagFieldMedia
% 11) MagFieldVacuum
% 12) MaxwellEqns
% 13) Radiation
% 14) Relativity
% 15) SteadyCurrent
% 16) SteadyEfield
% 17) Telegraph
% 18) Waveguides

\newcommand{\includechapter}[1]{\input{#1/#1}}
\newcommand{\includetikz}[2]{\input{#1/tikz/#2.tikz}}


\begin{document}

%% ==========================================================
%% Вставити після рядка 541 (після Attention "...дротах, так само, як вода тече по трубі.")
%% і перед \section{Скін-ефект} (рядок 544) у файлі MaxwellEqns.tex
%%
%% ПЛАН МАЛЮНКІВ (усі --- заглушки, tikz будуть додані пізніше):
%%   1. pic:FieldLinesTension  --- пружні нитки: натяг уздовж + бічне відштовхування
%%   2. pic:StressElement      --- довільний майданчик dS, нормаль n, сила dF = T n dS
%%   3. pic:StressDecoding     --- однорідне поле вздовж x: майданчик n=e_x (натяг) і n=e_y (тиск)
%%   4. pic:WireInField        --- провідник у зовнішньому однорідному полі (згущення/розрідження ліній)
%%   5. pic:PinchEffect        --- плазмовий шнур, колові лінії, тиск до осі
%%   6. pic:ChargedDroplet     --- заряджена крапля, радіальні лінії, тиск назовні
%% ==========================================================

%% --------------------------------------------------------
\subsection*{Натяги Максвелла: тиск і натяг силових ліній}
%% --------------------------------------------------------

Густина енергії поля $w$, яку ми щойно ввели, --- це не лише зручна бухгалтерська величина для теореми Пойнтінга. У гл.~\ref{ElEnergy} та
гл.~\ref{ElMagInduction} ми вже двічі, енергетичним методом, переконалися, що поле діє на заряди й струми з поверхневою силою, яка за модулем
дорівнює саме густині енергії: тиск, що розпирає обмотку соленоїда, і натяг, що стягує обкладки конденсатора, --- обидва чисельно рівні $w$.
Виявляється, що це не випадковий збіг для двох окремо взятих геометрій, а загальна властивість: \emphx{механічні натяги, що виникають у будь-якій
задачі про пондеромоторні сили в магнітному полі, завжди мають таку величину, ніби силові лінії поводяться як натягнуті пружні нитки}
(рис.~\ref{pic:FieldLinesTension}), які прагнуть скоротитися й розвивають уздовж себе натяг, а також відштовхують одна одну впоперек, створюючи
бічний тиск, --- причому і натяг, і тиск за модулем дорівнюють об'ємній густині енергії поля $\mu H^2/8\pi$. Цілком аналогічне твердження
справедливе й для електричного поля: яку б задачу про пондеромоторні сили ми не розглядали, результат завжди такий, ніби електричні силові лінії
мають поздовжній натяг і бічний тиск, кожен з яких дорівнює густині енергії $\epsilon E^2/8\pi$.

Саме уявлення про натяг уздовж і тиск упоперек силових ліній запровадили в науку Фарадей і Максвелл. Тепер, коли в нас є вираз для повної
густини енергії $w=\dfrac{\epsilon E^2}{8\pi}+\dfrac{\mu H^2}{8\pi}$ одразу для обох полів, можна надати цій ідеї точну математичну форму --- не
для однієї конкретної геометрії, а для площадки \emphx{довільної} орієнтації.

\begin{figure}[h!]\centering
%	\localinput{field-lines-tension.tikz}
%% --------------------------------------------------------
%% Рисунок-інтуїція: силові лінії як пружні нитки.
%% Панель (а) -- одна лінія прагне скоротитися (натяг уздовж).
%% Панель (б) -- дві сусідні паралельні лінії відштовхуються
%%              впоперек (бічний тиск).
%% --------------------------------------------------------
\begin{tikzpicture}[>=Latex, scale=1.15, every node/.style={font=\small}]

	% ================= Панель (а): натяг уздовж лінії =================
	% сама "пружна нитка" -- лінія у вигляді пружини
	\draw[very thick, decorate,
		decoration={coil, aspect=0.35, segment length=3.4pt, amplitude=1.7pt}]
		(-4.0,-1.3) -- (-4.0,1.3);

	% стрілки натягу: кінці лінії "тягнуться" одне до одного
	\draw[->, thick, blue!60!black] (-4.0,1.75) -- (-4.0,1.35);
	\draw[->, thick, blue!60!black] (-4.0,-1.75) -- (-4.0,-1.35);
	\node[blue!60!black, right] at (-3.85,1.55) {натяг};
	\node[blue!60!black, right] at (-3.85,-1.55) {натяг};

	\node at (-4.0,-2.35) {(а) натяг уздовж лінії};

	% ================= роздільник =================
	\draw[gray!70, dashed] (-1.6,-2.0) -- (-1.6,2.0);

	% ================= Панель (б): тиск упоперек ліній =================
	% дві сусідні паралельні "нитки"
	\draw[very thick, decorate,
		decoration={coil, aspect=0.35, segment length=3.4pt, amplitude=1.5pt}]
		(0.1,-1.3) -- (0.1,1.3);
	\draw[very thick, decorate,
		decoration={coil, aspect=0.35, segment length=3.4pt, amplitude=1.5pt}]
		(1.1,-1.3) -- (1.1,1.3);

	% стрілки бічного тиску: лінії штовхають одна одну впоперек
	\draw[->, thick, orange!85!black] (0.6,0.15) -- (0.15,0.15);
	\draw[->, thick, orange!85!black] (0.6,0.15) -- (1.05,0.15);
	\node[orange!85!black] at (0.6,0.5) {тиск};

	\node at (0.6,-2.35) {(б) тиск упоперек ліній};

\end{tikzpicture}

	\caption{Силові лінії поля як пружні нитки: уздовж лінії --- натяг (нитка прагне скоротитися), між сусідніми паралельними
		лініями --- бічний тиск (нитки взаємно відштовхуються)\label{pic:FieldLinesTension}}
\end{figure}

\begin{Attention}
	Варто одразу застерегти від буквального розуміння цієї картини: \emphx{жодних фізичних \enquote{ниток} у просторі не існує} --- силові лінії
	є суто геометричним образом, який ми вводимо для наочного зображення поля. Проте уявлення Фарадея й Максвелла залишається надзвичайно
	корисним інструментом: воно дає простий спосіб \emphx{визначити напрям і характер} механічних сил у електромагнітному полі, навіть коли пряме
	обчислення через силу Кулона чи силу Ампера було б громіздким.
\end{Attention}

\subsubsection*{Тензор натягів Максвелла}

Щоб описати силу на площадку довільної орієнтації, натяг і тиск зручно об'єднати в один об'єкт --- \emphx{тензор натягів Максвелла}. Для
площадки з одиничною нормаллю $\vect n$ сила, що діє на неї з боку поля (рис.~\ref{pic:StressElement}),
\begin{equation}\label{eq:StressDefinition}
	d\vect F = \hat{T}\,\vect n\, dS,
\end{equation}
де $\hat{T}$ --- тензор другого рангу; для електричного поля
\begin{equation}\label{eq:MaxwellStressTensor}
	\tcbhighmath{
		T_{ik} = \frac{\epsilon}{4\pi}\left(E_iE_k - \frac12\delta_{ik}E^2\right),
	}
\end{equation}
а для магнітного (заміною $\Efield\to\Hfield$, $\epsilon\to\mu$) ---
\begin{equation}\label{eq:MaxwellStressTensorB}
	\tcbhighmath{
		T_{ik} = \frac{1}{4\pi\mu}\left(H_iH_k - \frac12\delta_{ik}H^2\right).
	}
\end{equation}

\begin{figure}[h!]\centering
%	\localinput{stress-element.tikz}
%% --------------------------------------------------------
%% Рисунок до eq:StressDefinition: dF = \hat T n dS
%% Довільно орієнтований майданчик dS, нормаль n,
%% результуюча сила dF, розкладена на нормальну (dF_n)
%% і дотичну (dF_t) складові -- щоб показати, що dF
%% у загальному випадку НЕ колінеарна нормалі n.
%% --------------------------------------------------------
\begin{tikzpicture}[>=Latex, scale=1.5, every node/.style={font=\small}]

	% ---- однорідне поле: горизонтальні силові лінії (без лінії y=0, щоб не йшла крізь O) ----
	\foreach \y in {-1.8, -1.4, ...,1.8}{
		\draw[red!25, thick, -{Latex[length=2mm]}] (-2.4,\y) -- (2.1,\y);
	}
	\node[red!55] at (2.1,1.8) {$\Efield$};

	% ---- майданчик dS (похила пряма із штрихуванням) ----
	\coordinate (A) at (-0.376,-1.034);
	\coordinate (B) at ( 0.376, 1.034);
%	\draw[very thick] (A) -- (B);
%	\foreach \t in {0.08,0.22,...,0.95}{
%		\coordinate (P) at ($(A)!\t!(B)$);
%		\draw[thick] (P) -- ($(P)+(-0.20,-0.11)$);
%	}
%	\node at ($(B)+(0.10,0.20)$) {$dS$};

	\coordinate (O) at (0,0);
    \draw[black, fill=blue!50, opacity=0.5, rotate around={-25:(O)}] (O) circle(0.25 and 0.5);

	% ---- нормаль n (коротка, лише щоб показати напрям) ----
	\coordinate (N) at (-25:1);
	\draw[->, very thick] (O) -- (N);
	\node[above=1pt] at (N) {$\vect n$};

	% ---- результуюча сила dF (не колінеарна n!) ----
	\coordinate (F) at (-55:{2/cos((-55+25))});
	\draw[->, ultra thick, red!70!black] (O) -- (F)
		node[below right, red!70!black] {$d\vect F=\hat T\,\vect n\,dS$};

	% ---- розкладання dF на нормальну й дотичну складові ----
	\coordinate (Q) at (-25:2); % проєкція F на напрям n
	\draw[dashed] (O) -- (Q) node[pos=0.9, above] {$dF_n$};
	\draw[dashed] (Q) -- (F) node[midway, right=2pt] {$dF_t$};

	% ---- позначення точки прикладання ----
	\fill (O) circle (1.1pt);
\end{tikzpicture}
	\caption{Довільно орієнтований майданчик $dS$ з нормаллю $\vect n$ у полі: результуюча сила $d\vect F=\hat T\vect n\,dS$ в
		загальному випадку не колінеарна нормалі\label{pic:StressElement}}
\end{figure}

Переконаймося, що це узгоджується з уже відомим результатом (рис.~\ref{pic:StressDecoding}). Нехай $\Efield=E\vect e_x$. Тоді
з~\eqref{eq:MaxwellStressTensor}
\begin{equation*}
	T_{xx} = \frac{\epsilon E^2}{8\pi} = w_E, \qquad T_{yy}=T_{zz}=-\frac{\epsilon E^2}{8\pi}=-w_E,
\end{equation*}
а недіагональні компоненти дорівнюють нулю. Для майданчика $\vect n=\vect e_x$ (перпендикулярного полю) сила $d\vect F=T_{xx}\vect e_x\,dS=+w_E\vect e_x\,dS$
напрямлена вздовж поля --- це натяг, який тягне обкладки конденсатора (гл.~\ref{ElEnergy}). Для майданчика $\vect n=\vect e_y$ (паралельного
полю) сила $d\vect F=T_{yy}\vect e_y\,dS=-w_E\vect e_y\,dS$ напрямлена проти нормалі --- це тиск, який розпирає обмотку соленоїда
(гл.~\ref{ElMagInduction}). Обидва вже знайомі результати --- лише два окремі майданчики одного й того самого тензора, тепер \eqref{eq:MaxwellStressTensor}
одразу дає силу для будь-якої орієнтації площадки.

\begin{figure}[h!]\centering
%	\localinput{stress-decoding.tikz}
	\caption{Однорідне поле вздовж осі $x$: майданчик $\vect n=\vect e_x$ відчуває натяг (сила вздовж поля, всередину), а
		майданчик $\vect n=\vect e_y$ --- тиск (сила проти нормалі, назовні)\label{pic:StressDecoding}}
\end{figure}

\subsubsection*{Приклад: прямий провідник в однорідному магнітному полі}

Особливо наочно картина натягу й тиску працює тоді, коли треба лише \emphx{визначити напрям} сили, не обчислюючи її величину. Нехай прямий
провідник зі струмом $I$ вміщено в однорідне магнітне поле $\Hfield_0$, силові лінії якого --- паралельні прямі (рис.~\ref{pic:WireInField}).
Власне поле струму --- це, як відомо, концентричні кола навколо провідника. Накладемо обидві картини одна на одну: з одного боку від провідника
лінії власного поля течуть \emphx{в один бік} із лініями $\Hfield_0$, і сумарне поле там \emphx{згущується}; з іншого боку вони течуть
\emphx{назустріч} одна одній, частково компенсуючись, --- і сумарне поле там \emphx{розріджується}.

\begin{figure}[h!]\centering
%	\localinput{wire-in-uniform-field.tikz}
	\caption{Прямий провідник зі струмом в однорідному магнітному полі: складання власного поля струму (концентричні кола) з
		однорідним полем $\Hfield_0$ дає згущення ліній з одного боку провідника й розрідження --- з іншого\label{pic:WireInField}}
\end{figure}

Але паралельні силові лінії, як ми щойно бачили, взаємно \emphx{відштовхуються впоперек} --- отже там, де вони згущені, цей бічний тиск
найбільший. Провідник виявляється затиснутим між ділянкою підвищеного тиску (де лінії згущені) і ділянкою пониженого тиску (де вони розріджені),
а тому виштовхується в бік розрідження --- \emphx{перпендикулярно} і до провідника, і до початкового поля $\Hfield_0$. Це саме той напрям сили
Ампера $\vect F = \dfrac{I}{c}\left[\vect\ell\times\Hfield_0\right]$, який ми вже отримали раніше (гл.~\ref{MagFieldVacuum}) прямим
обчисленням, --- тепер же він випливає з самої лише картини силових ліній, без жодної формули.

\subsubsection*{Приклад: пінч-ефект у плазмі}

Ще один приклад, де вже недостатньо простого правила \enquote{тиск чи натяг} --- потрібна саме тензорна картина орієнтації ліній відносно
поверхні. Нехай циліндричний плазмовий шнур радіуса $a$ несе повний струм $I$ уздовж своєї осі. Магнітне поле $H(r)$, яке цей струм створює,
охоплює шнур концентричними колами --- на відміну від соленоїда, тут силові лінії \emphx{замкнені навколо} струму, а не паралельні йому
(рис.~\ref{pic:PinchEffect}). Уздовж кожної такої колової лінії поле чинить натяг (лінія \enquote{хоче} стягнутися в кільце меншого радіуса), а
це означає \emphx{тиск, спрямований до осі шнура}: поле стискає плазму, немов обруч.

\begin{figure}[h!]\centering
%	\localinput{pinch-effect.tikz}
	\caption{Пінч-ефект: колові силові лінії магнітного поля навколо струму в плазмовому шнурі створюють тиск, спрямований
		до осі шнура\label{pic:PinchEffect}}
\end{figure}

Це --- \emphx{пінч-ефект} (від англ.\ \emph{pinch}, \enquote{защемлення}): струмонесучий плазмовий шнур стискається власним магнітним полем.
Стиснення триває, доки магнітний тиск не зрівноважиться тепловим (кінетичним) тиском самої плазми:
\begin{equation}\label{eq:PinchBalance}
	\frac{H^2}{8\pi} \sim n k_{\text{B}} T,
\end{equation}
де $n$ --- концентрація частинок плазми, $T$ --- її температура. Саме пінч-ефект лежить в основі найпростіших схем магнітного утримання плазми в
термоядерних установках.

\subsubsection*{Приклад: заряджена крапля}

Натяг і тиск силових ліній можна побачити й на викривленій поверхні. Уявіть заряджену провідну краплю (наприклад, краплю ртуті чи мильну
бульбашку з надлишковим зарядом): силові лінії виходять з її поверхні перпендикулярно, назовні, в усіх напрямках (рис.~\ref{pic:ChargedDroplet}).
Як і в соленоїді, тут поле чинить \emphx{тиск упоперек ліній}, а лінії йдуть уздовж нормалі до поверхні --- отже цей тиск $p_E=\epsilon E^2/8\pi$
штовхає поверхню назовні, \enquote{роздуваючи} краплю.

\begin{figure}[h!]\centering
%	\localinput{charged-droplet.tikz}
	\caption{Заряджена крапля: радіальні силові лінії створюють тиск, спрямований назовні, якому протидіє поверхневий
		натяг рідини\label{pic:ChargedDroplet}}
\end{figure}

Цьому тиску протидіє поверхневий натяг рідини $\sigma$, що намагається стягнути краплю в кулю мінімальної площі. Поки заряд малий, поверхневий
натяг перемагає; але коли заряд перевищує критичне значення, електричний тиск перевищує поверхневий натяг, і крапля втрачає стійкість --- вона
витягується в конус і \enquote{вистрілює} тонкими зарядженими струменями. Це явище називають \emphx{нестійкістю Релея}, і воно ж лежить в основі
електроспрею й деяких методів мас-спектрометрії.

\begin{Attention}
	В усіх чотирьох прикладах --- прямому провіднику, пінч-ефекті, конденсаторі й соленоїді, зарядженій краплі --- ми жодного разу не
	підсумовували сили Кулона чи Ампера по окремих зарядах і струмах: досить було визначити, \emphx{як орієнтовані силові лінії відносно
	поверхні}. Це та сама механічна дія поля, що й потік енергії через вектор Пойнтінга, --- лише інша, \enquote{силова}, а не
	\enquote{енергетична}, проєкція одного й того самого факту: поле не просто переносить енергію, воно ще й чинить цілком реальну механічну дію
	на все, крізь що проходить. Коли пізніше ми дійдемо до електромагнітних хвиль, побачимо, що й тиск хвилі при її поглинанні речовиною --- не
	окремий ефект, а той самий натяг поля, який хвиля \enquote{несе} з собою.
\end{Attention}

\end{document}

